\section{An exact symmetry useful for comparing targets}
\label{sec:projective-cube-symmetry}

The first-layer bias has the same Gaussian variance as each input
coefficient. This gives a slightly larger symmetry than coordinate
permutations of the Boolean cube. The result below holds at finite width
and every prior scale, not only for the critical rate.

\begin{proposition}[Projective coordinate symmetry]
\label{prop:projective-cube-symmetry}
Write $y=(1,x_1,\ldots,x_d)$ and let $\pi$ be any permutation of
$\{0,\ldots,d\}$. Define
\[
 a_\pi(x)=y_{\pi(0)},\qquad
 T_\pi(x)_j=y_{\pi(j)}/y_{\pi(0)}\quad(1\le j\le d).
\]
The map $T_\pi$ is a bijection of the Boolean cube. If $z$ is a network
drawn from the Gaussian prior, then
\[
 x\longmapsto a_\pi(x)z(T_\pi(x))
\]
has exactly the same law as $z$. Consequently, writing
$f^\pi(x)=a_\pi(x)f(T_\pi(x))$, the target events $A_{N,f}$ and
$A_{N,f^\pi}$ have equal prior probabilities. At critical scaling,
$J_L(f)=J_L(f^\pi)$.

In particular, for every $m$ for which the indicated degrees exist,
all parity functions of degree $2m$ or $2m+1$ have the same prior
probability and the same critical rate. Constants and dictators are the
first such pair of degree classes.
\end{proposition}

\begin{proof}
Permute the $d+1$ columns of the first-layer matrix by $\pi$. Its
Gaussian law is unchanged. The permuted input vector satisfies
\[
 (y_{\pi(0)},\ldots,y_{\pi(d)})
       =a_\pi(x)(1,T_\pi(x)).
\]
Every subsequent hidden activation is odd, and there are no subsequent
biases. Thus the factor $a_\pi(x)\in\{-1,1\}$ passes through all hidden
layers and the linear readout. This proves the asserted identity in
law. The induced transformation of projective sign vectors is
invertible, which proves that $T_\pi$ is a bijection. The equality of
target-event probabilities follows by changing variables on the cube;
the rate equality follows either from these exact finite-width
identities or by the same transformation in the variational functional.

For the final assertion it suffices to swap coordinate zero with an
input coordinate $i$. If a parity has even degree and does not involve
$i$, the transformation adds $i$. If a parity has odd degree and does
involve $i$, it removes $i$. In the other two cases it leaves the parity
unchanged. Ordinary input permutations identify all subsets of a given
size.
\end{proof}

For uniform input sampling, transforming each example by $T_\pi$ and
multiplying its label by $a_\pi$ also identifies the corresponding
posterior learning problems, including Gibbs error and any fixed
posterior-confidence threshold. This exact symmetry does not identify
all parity degrees, and does not by itself give a sample-complexity
lower bound.
